\documentclass{book}
\usepackage[backend=biber,natbib=true,style=authoryear]{biblatex}
\addbibresource{/home/nguyen/1_NQBH/math/bib.bib}
\usepackage[utf8]{inputenc}
\usepackage{float}
\usepackage{graphicx}
\usepackage[colorlinks=true,linkcolor=blue,urlcolor=red,citecolor=magenta]{hyperref}
\usepackage{amsmath,amssymb,amsthm}
\allowdisplaybreaks
\numberwithin{equation}{section}
\newtheorem{assumption}{Assumption}[section]
\newtheorem{lemma}{Lemma}[section]
\newtheorem{corollary}{Corollary}[section]
\newtheorem{definition}{Definition}[section]
\newtheorem{proposition}{Proposition}[section]
\newtheorem{theorem}{Theorem}[section]
\newtheorem{notation}{Notation}[section]
\newtheorem{remark}{Remark}[section]
\newtheorem{example}{Example}[section]
\newtheorem{ques}{Question}[section]
\newtheorem{problem}{Problem}[section]
\newtheorem{conjecture}{Conjecture}[section]
\usepackage[left=0.5in,right=0.5in,top=1.5cm,bottom=1.5cm]{geometry}
\usepackage{fancyhdr}
\pagestyle{fancy}
\fancyhf{}
\addtolength{\headheight}{0pt}% obsolete
\lhead{\itshape \scriptsize \chaptername~\thechapter}
\rhead{\itshape \scriptsize \nouppercase{\leftmark}} %\nouppercase !
\renewcommand{\chaptermark}[1]{\markboth{#1}{}}
\cfoot{\thepage}

\title{\href{https://medium.com}{Medium}}
\author{Collector: Nguyen Quan Ba Hong}
\date{\today}

\begin{document}
\maketitle
\setcounter{secnumdepth}{6}
\setcounter{tocdepth}{6}
\tableofcontents

%------------------------------------------------------------------------------%

\part{\href{https://medium.com/@geniusturner360}{Genius Turner}}

\chapter{\href{https://medium.com/mind-cafe/einsteins-formula-for-a-happy-life-b29aff61a9c7}{Einstein's Formula for a Happy Life: There's a science to everything}}

Everything we do in life is for a reason.

And so, quite naturally, for ages humans have wondered: \textit{For what reason do I exist?}

%
Some men chase women, some women chase money, some of both chase vodka with the soft drink of choice, yet every such chase leads to the same pot at the end of the rainbow: happiness.

\begin{quotation}
    ``\textit{Happiness},'' \href{https://historyofeconomicthought.mcmaster.ca/aristotle/Ethics.pdf}{said Aristotle}, ``\textit{is the meaning and the purpose of life, the whole aim and end of human existence}.''
\end{quotation}
Armed with the above insight into the very reason for our existence, does not common-sense suggest it would be wise to inquire with the very man widely considered the smartest person in history?

%
If by chance you were to Google \textit{genius definition}, notice the following incredible inclusion:

\textbf{genius} [n]
\begin{enumerate}
    \item exceptional intellectual or creative power or other natural ability.
    
    ``she was a teacher of genius''
    
    \textit{synonyms}: \textbf{brilliance}, great intelligence, great intellect, great ability, \textbf{cleverness}, brains, \textbf{erudition, wisdom, sagacity}, find mind, \textbf{wit, artistry, flair}, creative power, precocity, precociousness;
    \item an exceptionally intelligent person or one with exceptional skill in a particular area of activity.
    
    ``a mathematical genius''
    
    \textit{synonyms}: brilliant person, mental giant, \textbf{mastermind}, Einstein, \textbf{intellectual, intellect, brain, highbrow, expert, master, artist, polymath}. 
\end{enumerate}
Doesn't seeing an actual person's last name - Einstein - listed along with abstract nouns stand out like seeing Shaquille O'Neal in a room filled with gymnasts?

%
In short, because writing is nothing but thinking on paper, to read the thoughts of another is to essentially converse with them.

And so, who better to ask about the secret to happiness than genius personified?

\section{Einstein's Formula for a Happy Life}
A few days before Einstein twirled into the Reaper's grim arms, his assistant - Dukas - found him in the hospital bed, ``in agony, unable to lift his head.''

%
Yet on the very next day, a mere 24 hours or so away from his death-day, Einstein ``asked Dukas to get him his glasses, papers, and pencil, and he proceeded to jot down a few calculations.''

%
``He worked as long as he could,'' \href{https://workbooks.colombo.ca/issue-37-15-april-2020/}{noted biographer} Walter Issacson, ``and when the pain got too great he went to sleep,'' for the final time.

Indeed, Einstein died doing the 1 thing he loved most - working.

Ah, circumstances reveal character!

%
``Genius is one percent inspiration,'' said Einstein, ``and 99 percent perspiration.''

Indeed, it's not by accident that no one has ever become great by accident.

After all, as \href{https://www.newyorker.com/magazine/1947/11/22/the-great-foreigner}{Einstein once noted}: ``Only a monomaniac gets what we commonly refer to as \textit{results}.''

%
Show me someone great and I'll show you someone obsessed.

Besides, what more is ``greatness'' than the child of an obsession?

%
For the above reason, when Einstein was asked for the secret to a happy life, though the questioner expected an answer long and sour, Einstein kept it short and sweet:
\begin{quotation}
    ``\textit{If you want to live a happy life, tie it to a goal, not to people or things}.''
\end{quotation}
Here lies Einstein's formula for a happy life.

\section{Work - the Ultimate Shelter from Life's Storms}
Given that great minds think alike for the same reason passengers boarded the same train of thought inevitably end up at the same destination, we should hardly be surprised that Newton and Einstein - arguably the 2 greatest scientists in history - when confronted with life's storms both used the same formula for a happy life.

%
In 1665, the \href{https://www.historic-uk.com/HistoryUK/HistoryofEngland/The-Great-Plague/}{bubonic plague} struck London.

Never before nor since has the world seen anything like the deadly pandemic. And just as a pandemic shut down today's stores and universities, the same held true back then.

History repeats itself indeed.

%
While the world's stage was seemingly crumbling underneath his feet, Isaac Newton applied the formula for a happy life.

How?

He merely tied his life to an abstract goal, not to physical people or concrete things.

Over the span of roughly 18 months, Newton would revolutionize science.

%
Newton was later asked how he discovered the law of gravity.

He replied:
\begin{quotation}
    ``\textit{By thinking about it all the time}.''
\end{quotation}
In short, because Newton had hitched his star to an abstract goal, which remains forever fixed, and not to life - which is forever unpredictable - he in effect found shelter from life's storms.

%
When Einstein's wife Elsa died, needless to say - he was crushed.

After all, according to Isaacson, Elsa served as a somewhat maternal figure to Einstein.

%
``She told him when to eat and where to go,'' \href{https://onlinereadfreenovel.com/walter-isaacson/page,171,42091-walter_isaacson_great_innovators_e-book_boxed_set.html}{Isaacson notes}.

``She packed his suitcases and doled out his pocket money. In public, she was protective of the man she called `the Professor.'''

%
Fortunately for Einstein, his formula for a happy life served as shelter from the storm.

\begin{quotation}
    ``\textit{As long as I am able to work, I must not and will not complain, because work is the only thing that gives substance to life}.''
\end{quotation}

\section{In Closing: Dream With Mind $+$ Chase with Body $=$ the Formula for a Happy Life}
When the holy man was asked does he feel good due to having given up all the worldly pleasures, he answered: ``It's not so much that I feel good, but rather I no longer feel bad.''

%
What Einstein fully grasped - regarding his concise formula for a happy life - appears to boil down to this:

\begin{quotation}
    \textit{``Happiness'' can be likened to a supermodel, all glammed up for the runway. Contentment, however, is that same model when she returns home at night and then removes all the makeup}.
\end{quotation}
Or to put it another way, ``He is richest,'' said Socrates, ``who is content with the least.''

%
Indeed, the best things in life are not only free but they're not even ``things.''

After all, never has a hand touched \textit{love}.

Never has an eye spotted \textit{peace}.

Never has a nose sniffed \textit{dreams}.

Here lies the DNA of Einstein's formula for a happy life.

%
From aardvark to zebra, we humans are the only creatures to have ever graced the world's stage that can ``dream.''

To picture a goal with our 3rd eye coupled with pouring our heart and soul into manifesting it before our 2 eyes is the very reason for this magical gift called a ``mind.''

%
The above insight may have best been summed up by a character in the classic play \textit{The Secret of Freedom}:

\begin{quotation}
    ``\textit{The only thing about a man that is man is his mind. Everything else you can find in a pig or a horse}.''
\end{quotation}
Perhaps there's much truth in the saying that we humans only value things we pay for.

%
Like those other abstract pleasures which came freely packaged at birth, such as \textit{love} and \textit{faith}, perhaps we take for granted our gift for dreaming up worthwhile goals, not to mention the ability to wed our lives to such dreaming.

%
Because example is better than precept, perhaps the occasion calls for ending this piece by recounting a real-life instance of how using Einstein's formula for a happy life sheltered 1 of my good friends from a brutal storm.

%
1 of my closest friends in college endured the ultimate heartbreak.

Her high school sweetheart and fiancé left her for another woman.

To say my pal was devastated would be an understatement.

Yet somehow$\ldots$ someway she embodied Maya Angelou's grand vision of a \href{https://www.poetryfoundation.org/poems/48985/phenomenal-woman}{Phenomenal Woman}.

%
How?

%
She simply rose from the canvas of doubt and fear, dusted herself off, and then poured all her time and energies into her studies.

A year or so later she got accepted into Loyola Law School.

%
Today, she's not only happily married but also a successful attorney.

In short, my pal used Einstein's formula for a happy life.

\section{The Takeaway}
Because Einstein was a mathematician at heart, he couldn't help but reduce the odds of achieving happiness down to an equation.

\begin{quotation}
    ``\textit{You have to learn the rules of the game},'' he concluded, ``\textit{and then you have to play better than anyone else}.''
\end{quotation}
Dear reader, so far as the Game of Life is concerned, it appears the rules are such as to demand that we - the only creatures armed with an imagination - use it specifically to envision an abstract goal, better known as a ``dream.''

%
Perhaps the above insight explains why Einstein went so far as to have called ``\textit{imagination more important than knowledge}.''

In short, though wedding your life to dream-chasing may not sound like the most glorious endeavor, as Einstein noted - such a course is nevertheless your safest bet.

%
Most of all, given that no mortal knows for certain what the next moment has in store - from triumph to tragedy - wisdom dictates treating your life as a boat and your goal as an anchor, which will, in turn, lend stability during life's storms.

%
In short, Einstein's formula for a happy life can best be summed as follows:

\begin{quotation}
    \textit{Find out what you do best; then, find out how you can pay most of your attention to doing it; then, get someone to pay you for having paid most of your attention to mastering that 1 thing}.
\end{quotation}

\section{Sources}
\begin{enumerate}
    \item Bartlett, Robert (2012). \textit{Aristotle's Nicomachean Ethics}.
    \item Isaacson, Walter (2007). \textit{Einstein: His Life and Universe}.
    \item Tucci, Niccolo (1947). \textit{The Great Foreigner}.\hfill$\square$
\end{enumerate}

%------------------------------------------------------------------------------%

\part{\href{https://medium.com/@ztrana}{Zat Rana}}

\chapter{\href{https://medium.com/personal-growth/the-most-important-skill-nobody-taught-you-9b162377ab77}{The Most Important Skill Nobody Taught You}}

Before dying at the age of 39, Blaise Pascal made huge contributions to both physics and mathematics, notably in fluids, geometry, and probability.

%
This work, however, would influence more than just the realm of the natural sciences.

Many fields that we now classify under the heading of social science did, in fact, also grow out of the foundation he helped lay.

%
Interestingly enough, much of this was done in his teen years, with some of it coming in his twenties.

As an adult, inspired by a religious experience, he actually started to move towards philosophy and theology.

%
Right before his death, he was hashing out fragments of private thoughts that would later be released as a collection by the name of \href{https://designluck.com/recommends/pascals-pensees/}{Pensées}.

%
While the book is mostly a mathematician's case for choosing a life of faith and belief, the more curious thing about it is its clear and lucid ruminations on what it means to be human.

It's a blueprint of our psychology long before psychology was deemed a formal discipline.

%
There is enough thought-provoking material in it to quote, and it attacks human nature from a variety of different angles, but one of its most famous thoughts aptly sums up the core of his argument:

\begin{quotation}
    ``\textit{All of humanity's problems stem from man's inability to sit quietly in a room alone}.''
\end{quotation}
According to Pascal, we fear the silence of existence, we dread boredom and instead choose aimless distraction, and we can't help but run from the problems of our emotions into the false comforts of the mind.

%
The issue at the root, essentially, is that we never learn the art of solitude.

\section{The Perils of Being Connected}
Today, more than ever, Pascal's message rings true.

If there is one word to describe the progress made in the last 100 years, it's connectedness.

%
Information technologies have dominated our cultural direction.

From the telephone to the radio to the TV to the internet, we have found ways to bring us all closer together, enabling constant worldly access.

%
I can sit in my office in Canada and transport myself to practically anywhere I want through Skype.

I can be on the other side of the world and still know what is going on at home with a quick browse.

%
I don't think I need to highlight the benefits of all this.

But the downsides are also beginning to show.

Beyond the current talk about privacy and data collection, there is perhaps an even more detrimental side-effect here.

%
\textit{We now live in a world where we're connected to everything except ourselves}.

%
If Pascal's observation about our inability to sit quietly in a room by ourselves is true of the human condition in general, then the issue has certainly been augmented by an order of magnitude due to the options available today.

%
The logic is, of course, seductive.

Why be alone when you never have to?

%
Well, the answer is that never being alone is not the same thing as never feeling alone.

Worse yet, the less comfortable you are with solitude, the more likely it is that you won't know yourself.

And then, you'll spend even more time avoiding it to focus elsewhere.

In the process, you'll become addicted to the same technologies that were meant to set you free.

%
Just because we can use the noise of the world to block out the discomfort of dealing with ourselves doesn't mean that this discomfort goes away.

%
Almost everybody thinks of themselves as self-aware.

They think they know how they feel and what they want and what their problems are. But the truth is that very few people really do.

And those that do will be the 1st to tell how fickle self-awareness is and how much alone time it takes to get there.

%
In today's world, people can go their whole lives without truly digging beyond the surface-level masks they wear; in fact, many do.

%
We are increasingly out of touch with who we are, and that's a problem.

\section{Boredom as a Mode of Stimulation}
If we take it back to the fundamentals - and this is something Pascal touches on, too - our aversion to solitude is really an aversion to boredom.

%
At its core, it's not necessarily that we are addicted to a TV set because there is something uniquely satisfying about it, just like we are not addicted to most stimulants because the benefits outweigh the downsides. Rather, what we are really addicted to is a state of not-being-bored.

%
Almost anything else that controls our life in an unhealthy way finds its root in our realization that we dread the nothingness of nothing. We can't imagine just \textit{being} rather than \textit{doing}.

And therefore, we look for entertainment, we seek company, and if those fail, we chase even higher highs.

%
We ignore the fact that never facing this nothingness is the same as never facing ourselves.

And never facing ourselves is why we feel lonely and anxious in spite of being so intimately connected to everything else around us.

%
Fortunately, there is a solution.

The only way to avoid being ruined by this fear - like any fear - is to face it.

It's to let the boredom take you where it wants so you can deal with whatever it is that is really going on with your sense of self.

That's when you'll hear yourself think, and that's when you'll learn to engage the parts of you that are masked by distraction.

%
The beauty of this is that, once you cross that initial barrier, you realize that being alone isn't so bad.

Boredom can provide its own stimulation.

%
When you surround yourself with moments of solitude and stillness, you become intimately familiar with your environment in a way that forced stimulation doesn't allow.

The world becomes richer, the layers start to peel back, and you see things for what they really are, in all their wholeness, in all their contradictions, and in all their unfamiliarity.

%
You learn that there are other things you are capable of paying attention to than just what makes the most noise on the surface.

Just because a quiet room doesn't scream with excitement like the idea of immersing yourself in a movie or a TV show doesn't mean that there isn't depth to explore there.

%
Sometimes, the direction that this solitude leads you in can be unpleasant, especially when it comes to introspection - your thoughts and your feelings, your doubts and your hopes - but in the long-term, it's far more pleasant than running away from it all without even realizing that you are.

%
Embracing boredom allows you to discover novelty in things you didn't know were novel; it's like being an unconditioned child seeing the world for the first time.

It also resolves the majority of internal conflicts.

\section{The Takeaway}
The more the world advances, the more stimulation it will provide as an incentive for us to get outside of our own mind to engage with it.

%
While Pascal's generalization that a lack of comfort with solitude is the root of all our problems may be an exaggeration, it isn't an entirely unmerited one.

%
Everything that has done so much to connect us has simultaneously isolated us.

We are so busy being distracted that we are forgetting to tend to ourselves, which is consequently making us feel more and more alone.

%
Interestingly, the main culprit isn't our obsession with any particular worldly stimulation.

It's the fear of nothingness - our addiction to a state of not-being-bored.

We have an instinctive aversion to simply \textit{being}.

%
Without realizing the value of solitude, we are overlooking the fact that, once the fear of boredom is faced, it can actually provide its own stimulation.

And the only way to face it is to make time, whether every day or every week, to just sit - with our thoughts, our feelings, with a moment of stillness.

%
The oldest philosophical wisdom in the world has one piece of advice for us: \textit{know yourself}.

And there is a good reason why that is.

%
Without knowing ourselves, it's almost impossible to find a healthy way to interact with the world around us.

Without taking time to figure it out, we don't have a foundation to built the rest of our lives on.

%
Being alone and connecting inwardly is a skill nobody ever teaches us.

That's ironic because it's more important than most of the ones they do.

%
\textit{Solitude may not be the solution to everything, but it certainly is a start}.\hfill$\square$

%------------------------------------------------------------------------------%



%------------------------------------------------------------------------------%

\printbibliography[heading=bibintoc]
\end{document}